\documentclass[11pt,a4paper]{article}
\usepackage{kotex}
\usepackage{fontspec}
\setmainfont{Times New Roman}
\setmainhangulfont{AppleMyungjo}
\setsanshangulfont{Apple SD Gothic Neo}
\usepackage[a4paper,top=18mm,bottom=18mm,left=20mm,right=20mm]{geometry}
\usepackage{fancyhdr}
\setlength{\headheight}{24pt}
\usepackage{enumitem}
\usepackage{mdframed}
\usepackage{booktabs}
\linespread{1.18}
\setlength{\parindent}{0pt}
\pagestyle{fancy}
\fancyhf{}
\renewcommand{\headrulewidth}{0.6pt}
\fancyhead[L]{\sffamily\bfseries 2026 고1 영어 변형 - 정답 및 해설 지문 37}
\fancyhead[R]{\sffamily 작업기억과 장기기억}
\fancyfoot[C]{\framebox[16mm][c]{\thepage}}
\newcommand{\qno}[1]{\par\vspace{10pt}{\sffamily\bfseries #1.}\ }
\newmdenv[linewidth=0.6pt,roundcorner=3pt,innertopmargin=7pt,innerbottommargin=7pt,
          innerleftmargin=9pt,innerrightmargin=9pt,skipabove=5pt,skipbelow=5pt]{pbox}
\begin{document}
\begin{center}
{\fontsize{15}{17}\selectfont\sffamily\bfseries [지문 37] 정답 및 해설}
\end{center}
\vspace{6pt}
\begin{center}
\renewcommand{\arraystretch}{1.35}
\begin{tabular}{cccc}
\toprule
문항 & 유형 & 정답 & 배점 \\
\midrule
1 & 글 순서 배열 & ④ & 3점 \\
2 & 빈칸추론 & ③ & 2점 \\
3 & 서술형 해석 & 모범답안 참조 & 2점 \\
4 & 요약문 & ① & 2점 \\
\bottomrule
\end{tabular}
\end{center}

\qno{1}\textbf{글 순서 배열 - 정답: ④ (C)-(A)-(B)}
\begin{pbox}
주어진 글은 작업기억은 의식적으로 인식되지만 장기기억은 직접 접근할 수 없다는 대비를 제시한다. (C)는 이 대비를 이어 받아 장기기억의 정보를 작업기억으로 가져와야 접근할 수 있다고 설명하고, 배우 예시를 시작한다. (A)는 그 예시를 확장하여 작업기억의 용량 한계 때문에 모든 배우를 동시에 의식할 수 없다고 말한다. (B)는 그러나 하나씩 떠올려 말할 때는 작업기억이 장기기억에서 정보를 검색한다는 결론을 제시한다.
\end{pbox}

\qno{2}\textbf{빈칸추론 - 정답: ③}
\begin{pbox}
빈칸은 ``Try and hold all the actors you know in your explicit awareness''의 \textit{explicit awareness}를 바꾼 자리이다. 글의 핵심 대비상 의식적으로 붙잡고 있는 기억은 작업기억이므로 ③ \textit{conscious working memory}가 정답이다.

\medskip
① long-term storage는 직접 의식에 올라와 있지 않은 저장소라서 반대이고, ④ subconscious recall system도 ``explicit''과 반대이다. ②와 ⑤는 본문의 기억 체계 설명과 관련이 약하다.
\end{pbox}

\qno{3}\textbf{서술형 해석}
\begin{pbox}
\textbf{대상 문장}: \textit{Your working memory can access your long-term memory and bring those memories to awareness.}

\medskip
\textbf{모범답안}: 당신의 작업기억은 장기기억에 접근하여 그 기억들을 의식으로 불러올 수 있다.

\medskip
\textbf{채점 기준}: ``working memory''를 작업기억으로, ``access long-term memory''를 장기기억에 접근하다로, ``bring those memories to awareness''를 그 기억들을 의식으로 불러오다/가져오다로 옮기면 정답.
\end{pbox}
\vspace{8pt}
\qno{4}\textbf{요약문 - 정답: ①}
\begin{pbox}
장기기억의 정보는 직접 의식되지 않으며, 작업기억이 그것을 하나씩 \textit{retrieves}(검색/인출)할 때 의식에 올라온다. 또한 작업기억은 한 번에 담을 수 있는 양이 제한되어 있으므로 (B)는 \textit{limited}가 적절하다.
\end{pbox}

\end{document}
